\chapter{Výsledky studentské práce}
\section{Návrh jednotlivých prvků chytré domáctnosti}
\subsection{Mikrokontroler}
%Mikrokontroler Xiao ESP32-C6 od společnosti Seeed Studio byl vybrán pro všechny tři prvky chytré domácnosti. Tento mikrokontroler je vybaven čipem ESP32-C6, jež disponuje rádiem IEEE 802.0.11ax pro Wi-Fi 6, IEEE 802.15.1 pro Bluetooth 5.3 a IEEE 802.15.4 pro Thread a podporuje protokol Matter. Je vybaven 4 MB flash pamětí a 400 kB SRAM pamětí. Pro napájení je možné použít USB-C konektor nebo pájecí plošky pro připojení baterie. Samotný mikrokontroler je vybaven nabíjecím obvodem pro Li-ion/Li-Po baterie. Mikrokontroler je rovněž vybaven konektorem MHF4 pro připojení externí antény. 
Pro všechny tři prvky chytré domácnosti byl zvolen vývojový modul Xiao ESP32-C6 od společnosti Seeed Studio~\cite{seeed-xiao-wiki}. Jeho jádrem je čip ESP32-C6, který nabízí rádio IEEE~802.11ax pro Wi-Fi~6, IEEE~802.15.1 pro Bluetooth~5.3 a IEEE~802.15.4 pro Thread, a podporuje tak protokol Matter~\cite{espressif-esp32c6-ds}. Modul disponuje 4~MB flash paměti a 400~kB SRAM. Napájet jej lze přes konektor USB-C nebo přes pájecí plošky určené pro připojení baterie, přičemž nabíjení Li-ion/Li-Po článků zajišťuje přímo integrovaný nabíjecí obvod. Pro externí anténu je k dispozici konektor MHF4~\cite{seeed-xiao-wiki,seeed-xiao-product}.
\subsection{Čidlo teploty}
\subsubsection{Výběr teplotního čidla}
%Bylo porovnáno několik běžně dostupných teplotních čidel z hlediska přesnosti měření, rychlosti odezvy a ceny. Jako referenční bylo použito čidlo SEN63C, které je součástí komerčního chytrého senzoru IKEA ALPSTUGA~\cite{sensirion-alpstuga,ikea-alpstuga,sensirion-sen6x}. Srovnávána byla čidla DHT11~\cite{aosong-dht11}, BMP180~\cite{bosch-bmp180} a IKEA TIMMERFLOTTE~\cite{ikea-timmerflotte}.
Bylo porovnáno několik běžně dostupných čidel teploty a relativní vlhkosti z hlediska systematické odchylky, rozptylu a lineárnosti odezvy. Jako referenční bylo použito čidlo SHT40~\cite{sensirion-sht4x}, vůči kterému byla vyhodnocena čidla SEN63C, které je součástí komerčního chytrého senzoru IKEA ALPSTUGA~\cite{sensirion-sen6x,sensirion-alpstuga,ikea-alpstuga}, dále DHT11~\cite{aosong-dht11}, BMP180~\cite{bosch-bmp180} a IKEA TIMMERFLOTTE~\cite{ikea-timmerflotte}. Čidlo BMP180 bylo hodnoceno pouze z hlediska teploty.
\begin{table}[ht]
\centering
\small
\setlength{\tabcolsep}{4pt}
\caption{Porovnání čidel (teplota), rozdíl čidla a referenčního SHT40 (jen klidné minuty)}
\label{tab:porovnani-teplota_klid}
\resizebox{\textwidth}{!}{%
\begin{tabular}{lrrrrrrr}
\hline
Čidlo & $n$ & $\bar{d}$ [$^\circ$C] & $s_d$ [$^\circ$C] & LoA [$^\circ$C] & RMSE [$^\circ$C] & $|d|\le 0{,}7$ [\%] & směrnice \\
\hline
ALPSTUGA (SEN63C) & 1916 & $0{,}43$ & $0{,}07$ & $0{,}29$ až $0{,}57$ & $0{,}44$ & $100{,}0$ & $1{,}16$ \\
BMP180 & 1916 & $-0{,}75$ & $0{,}02$ & $-0{,}80$ až $-0{,}71$ & $0{,}75$ & $0{,}5$ & $1{,}04$ \\
DHT11 & 1916 & $0{,}00$ & $0{,}04$ & $-0{,}09$ až $0{,}08$ & $0{,}04$ & $100{,}0$ & $0{,}98$ \\
TIMMERFLOTTE & 1916 & $-0{,}43$ & $0{,}08$ & $-0{,}58$ až $-0{,}28$ & $0{,}44$ & $100{,}0$ & $0{,}87$ \\
\hline
\end{tabular}%
}
\vspace{2mm}

\begin{minipage}{0.95\textwidth}
\footnotesize
\raggedright
\textbf{Vysvětlivky:} Rozdíl $d_i$ je hodnota čidla minus hodnota referenčního
čidla SHT40. $n$ -- počet párových měření; $\bar{d}$ -- střední hodnota
rozdílů; $s_d$ -- směrodatná odchylka rozdílů; LoA -- meze shody
$\bar{d} \pm 1{,}96\,s_d$~\cite{bland-altman-1986}; RMSE -- střední
kvadratická chyba; $|d|\le 0{,}7$ -- podíl měření, u nichž absolutní hodnota
rozdílu nepřekročila 0,7~$^\circ$C; směrnice -- sklon regresní přímky
(hodnota čidla jako funkce hodnoty SHT40).
\end{minipage}
\end{table}

\begin{table}[ht]
\centering
\small
\setlength{\tabcolsep}{4pt}
\caption{Porovnání čidel (vlhkost), rozdíl čidla a referenčního SHT40 (jen klidné minuty)}
\label{tab:porovnani-vlhkost_klid}
\resizebox{\textwidth}{!}{%
\begin{tabular}{lrrrrrrr}
\hline
Čidlo & $n$ & $\bar{d}$ [\%\,RH] & $s_d$ [\%\,RH] & LoA [\%\,RH] & RMSE [\%\,RH] & $|d|\le 3{,}0$ [\%] & směrnice \\
\hline
ALPSTUGA (SEN63C) & 1690 & $-7{,}53$ & $0{,}45$ & $-8{,}42$ až $-6{,}65$ & $7{,}55$ & $0{,}0$ & $1{,}04$ \\
DHT11 & 1690 & $-0{,}22$ & $0{,}55$ & $-1{,}31$ až $0{,}87$ & $0{,}60$ & $100{,}0$ & $1{,}15$ \\
TIMMERFLOTTE & 1690 & $-4{,}88$ & $0{,}25$ & $-5{,}36$ až $-4{,}40$ & $4{,}88$ & $0{,}0$ & $0{,}94$ \\
\hline
\end{tabular}%
}
\vspace{2mm}

\begin{minipage}{0.95\textwidth}
\footnotesize
\raggedright
\textbf{Vysvětlivky:} Význam veličin $n$, $\bar{d}$, $s_d$, LoA, RMSE a směrnice
je uveden pod tabulkou~\ref{tab:porovnani-teplota_klid}. Rozdíl $d_i$ je
hodnota čidla minus hodnota referenčního čidla SHT40. $|d|\le 3{,}0$ -- podíl
měření, u nichž absolutní hodnota rozdílu nepřekročila 3,0~\%\,RH.
\end{minipage}
\end{table}
\subsection{Chytré světlo}
Pro návrh chytrého světla bylo použito adresovatelné LED světlo WS2812B, které je běžně dostupné na trhu. Toto světlo umožňuje ovládání jednotlivých LED diod. 
\subsection{Bateriový vypínač}

\section{Návrh scémat}
\subsection{Návrh čidla teploty}
\subsection{Návrh chytrého světla}
\subsection{Návrh vypínače}
